% !PRINTED 561
\weberpage{561}
%<<CONT
dinates of a point $\xi$ in a space of $n$ dimensions
$R_\xi$, and call the expression

\weq{7}{\varrho = \sqrt{(\xi_1' - \xi_1'')^2 + (\xi_2' - \xi_2'')^2 + \cdots + (\xi_n' - \xi_n'')^2},}

the \gesperrt{distance} of the two points $\xi'$, $\xi''$.

If in (3) we let the variables $x_i$ run through all integer
values, the associated values of the $\xi_i$ give us a system of
points in the space $R_\xi$. The point that corresponds to the
values $x_i = 0$ is the initial point of the coordinates, or the
\gesperrt{origin}. The totality of these points remains unchanged if we
replace $x_i$ by $x_i + x_i^{(0)}$, where $x_1^{(0)}, x_2^{(0)},
\ldots, x_n^{(0)}$ denotes a fixed system of integers. We can also
express this fact by saying that the system of points of the
$\xi_i$ comes to cover itself when it undergoes a parallel
displacement by which the origin arrives at the place of an
arbitrarily given other point of the system; or again: The points
of the system are situated with respect to each of its points in
the same way as with respect to every other.

We shall call this system a \gesperrt{point lattice} and its
individual points the \gesperrt{lattice points}.

By (2) and (7) the value of $f$ is the square of the distance of
the point $\xi$ from the origin, and consequently $\sqrt{M}$ is
the \gesperrt{smallest distance between two lattice points} that occurs in
the whole system.

Besides the distances in the space $R_n$ we have also to consider
volumes, i.\,e. the values of the $n$-fold integral

\weq{8}{V_\xi = \int\!\!\int \ldots \int d\xi_1\, d\xi_2 \ldots d\xi_n,}

in which the variables $\xi_1, \xi_2, \ldots, \xi_n$ fill some
finite region $G_\xi$, which is determined, say, by one or more
inequalities. Thus the region bounded by the condition

\weq{9}{\xi_1^2 + \xi_2^2 + \cdots + \xi_n^2 \leqq 1 \text{ or also } \Sigma(\xi_i - \xi_i^0)^2 \leqq 1}

for arbitrary values $\xi_i^0$ is an \gesperrt{$n$-dimensional sphere} of
radius 1, whose volume we shall denote by $K_n$. From every
region $G_\xi$ we can derive a \gesperrt{similar region} $G_\xi^{(t)}$ by
enlarging the coordinates all in the ratio $t^{\frac{1}{n}} : 1$.
If $V^{(t)}$ is the volume of the new region, then by (8)

\weq{10}{V^{(t)} = t\,V.}
